\section{Discussion}

The framework suggests several structural conclusions. First, recognition should be modeled as a policy variable rather than a passive label: recognition can change the mapping from primitive randomness to trajectories. Second, ordinary policy adaptation and first-person conditioning are distinct but interacting mechanisms. Third, branch-wide recognition coherence is not the same quantity as single-observer marginal uplift.

This distinction is especially important for adaptive agents. A better ordinary policy can rescue states that an accessibility selector would otherwise downweight, producing a negative interaction while total first-person value still rises. In this sense, ordinary adaptation and QBS conditioning can be partial substitutes without being individually harmful.

The predictive-alignment analysis adds a separate structural point. In the clean score-measurable case, favorable ordering of conditional expected outcome and accessibility is enough to determine the covariance sign. In the more general model, accessibility can contain variation not explained by the predictive score, and the exact total-covariance decomposition isolates that residual dependence. This makes explicit what a richer observer model must supply: not only an expected accessibility response to the agent's information state, but also a defensible account of residual outcome--accessibility dependence.

The explained-variance form of the residual certificate is useful as an interpretation device rather than a new physical principle. It separates three quantities: how much outcome variation the signal explains, how much accessibility variation it explains, and whether the two explained components move in the same direction. The resulting condition is deliberately conservative because it protects against worst-case residual anti-correlation.

A further implication concerns information. Recognition can have nonnegative option value when it expands the feasible policy set and old policies remain available. This value-of-information statement is logically distinct from the optimality of any particular recognition-dependent policy.

A major identifiability boundary is classical. When accessibility is bounded, the normalized first-person law is exactly the conditional law obtained by recording each base trajectory with probability proportional to its accessibility. For any integrable nonnegative accessibility weight, the same law is also exactly representable as a classical size-biased record distribution in which the expected record multiplicity is proportional to accessibility; bounded truncations approximate the general law in total variation. Thus observer-conditioned outcome data that determine only the weighted law cannot by themselves distinguish an Everett accessibility mechanism from a behavior-matched classical ascertainment mechanism. This does not make the physical mechanisms equivalent, but it raises the evidential burden: a specifically Everettian bridge must be derived from independent physics or make additional observable, interventional, or sequential predictions not shared by the classical selection null.

The same problem persists across multiple policies or experimental contexts if the competing classical null is allowed to retune its selection channel independently in every context. For each absolutely continuous target observer law, a context-specific record-size-bias construction can reproduce that law exactly. Hence repeated interventions are mechanism-identifying only after both sides are constrained by predeclared cross-context structure. A simple restricted null with one shared selection function implies proportional Radon--Nikodym density ratios across contexts on a common support, but turning that mathematical restriction into an experiment additionally requires the base and selected laws to be identifiable on an aligned observable state space. If the restriction exists only on a latent branch description and can disappear under projection to actual records, it does not by itself provide an operational bridge test.

The main novelty claim should therefore be framed structurally: recognition-dependent policy changes jointly determine trajectories and observer-indexed accessibility; predictive information can organize their alignment; and the resulting ordinary-policy, conditioning, residual-dependence, and cross-context identification burdens can be made explicit and statistically audited. The standard probability identities, ascertainment representations, and unrestricted context-specific selection constructions used in those analyses are not themselves the novelty claim. Any stronger physical claim depends on an Everett bridge that supplies independently motivated shared structure and survives comparison with classical selection alternatives.
